% Source: https://artofproblemsolving.com/community/c3h1194825_pascals_triangle_marathon
% Stable identity: aops-topic-1194825-q01
\textbf{(a)} For $n\ge k+1$,
\[
\dfrac{\binom n{k+1}}{\binom nk}=\dfrac{\dfrac{n!}{(k+1)!\,(n-k-1)!}}{\dfrac{n!}{k!\,(n-k)!}}=\dfrac{k!\,(n-k)!}{(k+1)!\,(n-k-1)!}=\dfrac{n-k}{k+1}.
\]
This is a linear function of $n$, so consecutive ratios differ by the constant $\dfrac{1}{k+1}$: the ratios $\dfrac{1}{k+1},\dfrac{2}{k+1},\dfrac{3}{k+1},\dots$ form an arithmetic sequence with common difference $\boxed{\dfrac{1}{k+1}}$. For example, diagonal $2$ over diagonal $1$ gives $\dfrac{1}{2},\dfrac{3}{3},\dfrac{6}{4},\dfrac{10}{5},\dots=\dfrac{1}{2},1,\dfrac{3}{2},2,\dots$

\textbf{(b)} For diagonals $3$ and $1$,
\[
\dfrac{\binom n3}{\binom n1}=\dfrac{(n-1)(n-2)}{6},
\]
which is quadratic in $n$. For $n=3,4,5,6$ the ratios are $\dfrac{1}{3},1,2,\dfrac{10}{3}$, whose differences $\dfrac{2}{3},1,\dfrac{4}{3}$ are not constant. So the ratios of two non-adjacent diagonals need not form an arithmetic sequence, and the claim holds only for adjacent diagonals (in general the ratio $\dfrac{\binom nj}{\binom nk}$ for $j>k$ is a polynomial in $n$ of degree $j-k$).
